% Draft manuscript for the VG-SAE project.
% Compile from this directory with: latexmk -pdf apssamp.tex
\documentclass[aps,pre,reprint,superscriptaddress,amsmath,amssymb,nofootinbib]{revtex4-2}

\usepackage{bm}
\usepackage{hyperref}
\hypersetup{hidelinks}

\newcommand{\E}{\mathbb{E}}

\begin{document}

\title{\texorpdfstring{$L_0$}{L0} Is Not a Free Parameter: A Statistical-Mechanical Variational Garrote for Sparse Autoencoders}

\author{Yeongwoo Song}
\affiliation{Department of Physics, Korea Advanced Institute of Science and Technology, Daejeon 34141, Korea}

\author{Hawoong Jeong}
\email{hjeong@kaist.edu}
\affiliation{Department of Physics, Korea Advanced Institute of Science and Technology, Daejeon 34141, Korea}
\affiliation{Center for Complex Systems, Korea Advanced Institute of Science and Technology, Daejeon 34141, Korea}

\author{Junghyo Jo}
\email{jojunghyo@snu.ac.kr}
\affiliation{Department of Physics Education, Seoul National University, Seoul 08826, Korea}
\affiliation{Center for Theoretical Physics and Artificial Intelligence Institute, Seoul National University, Seoul 08826, Korea}
\affiliation{School of Computational Sciences, Korea Institute for Advanced Study, Seoul 02455, Korea}

\date{\today}

\begin{abstract}
Sparse autoencoders (SAEs) are usually selected on a sparsity--reconstruction
frontier, treating the mean number of active latents, $L_0$, as an operating
point.  This is insufficient for feature recovery: at an incorrectly low
$L_0$, an SAE can mix correlated features and reconstruct better than a model
constrained to the ground-truth directions at the same sparsity.  We formulate
an SAE as a statistical-mechanical system with input-dependent binary
occupations.  A mean-field posterior over these occupations turns the Gaussian
reconstruction energy, its Bernoulli
fluctuation correction, the support prior, and the posterior entropy into a
Variational Garrote free energy.  Its posterior expectation defines a soft
occupation conditional on the sparsity field, rather than a prescribed Top-$K$
budget.  We propose a three-stage program to test whether density-sweep
transitions locate the feature-recovery regime on a controlled sparse-coding
model, the official SynthSAEBench generator, and real language-model
activations, including settings where reconstruction alone can favor an
incorrect representation.
\end{abstract}

\maketitle

\section{Introduction}
\label{sec:introduction}

Sparse autoencoders seek a small set of latent directions whose activations
explain a high-dimensional representation.  In current practice, sparsity is
often imposed by fixing $K$ or by tuning a penalty and then comparing models
on a sparsity--reconstruction curve.  This procedure implicitly treats the
mean number of active latents, $L_0$, as a free design parameter: lower
reconstruction error at a chosen $L_0$ is taken as evidence for a better SAE\@.

That inference can fail.  Chanin and Garriga-Alonso showed that when the SAE
$L_0$ lies below the activity of the generating features, reconstruction loss
can favor decoder directions that mix correlated
features~\cite{chanin2026sparse}.  In their controlled examples, the mixed
dictionary can reconstruct better than a model constrained to the ground-truth
feature directions at the same incorrect $L_0$.  The same work finds
degenerate feature mixtures at excessive $L_0$
empirically~\cite{chanin2026sparse}.  Consequently,
reconstruction fidelity is necessary but not sufficient for feature recovery,
and a Pareto frontier cannot by itself determine the correct sparse
representation.  For a learned dictionary with $L$ decoder atoms, their
statistic
\begin{equation}
 c_{\mathrm{dec}}
 = \binom{L}{2}^{-1}\sum_{1\leq j<j'\leq L}
   \left|\cos(\bm d_j,\bm d_{j'})\right|
 \label{eq:cdec}
\end{equation}
is minimized near the generating $L_0$ in toy models.  On language-model SAEs,
its low-$L_0$ elbow coincides with peak sparse-probe performance, making it a
useful but non-identifying diagnostic.

We take the corresponding modeling position: for a generative code
$\bm z^\star_\mu$ with specified feature model and granularity,
\begin{equation}
 L_0^\star = \E\!\left[\lVert \bm z^\star_\mu\rVert_0\right]
\end{equation}
is an estimand, not an arbitrary sparsity budget.  This statement is
necessarily model dependent.  Changing the dictionary width or the feature
granularity changes the relevant occupation number, and real activations do
not provide an observed ground-truth $L_0^\star$.

The Variational Garrote (VG) supplies a statistical-mechanical language for
this problem.  VG introduces explicit binary selection variables, replaces
their intractable posterior by a mean-field distribution, and optimizes a
variational free energy rather than a reconstruction term
alone~\cite{soh2025variational}.  Its mask values are order parameters, its
precision is analogous to inverse temperature, and its sparsity cost is an
external field coupled to occupation.  Across the densities studied by Soh
\emph{et al.}, cross-realization selection uncertainty is smallest near matched
model and data density and rises under under- or over-selection; they fit the
full curve to estimate the underlying density~\cite{soh2025variational}.

Here we extend that construction from global regression-variable selection to
input-dependent SAE activity.  We do not remove reconstruction from the
objective; we embed it as a likelihood energy inside a normalized free energy
that also accounts for support fluctuations, entropy, and a sparse prior.  The
training target is therefore the variational free-energy surrogate, not bare
reconstruction energy.  The resulting VG-SAE does not set $L_0$ directly: its
fitted occupation is an output order parameter conditional on the sparsity
field.  We define a three-stage program to test whether occupation and
uncertainty transitions co-localize with feature-recovery or downstream optima.
Because $\gamma$ remains external, this is a sweep-based estimator, not
parameter-free identification.

\section{Variational Garrote SAE}
\label{sec:method}

Let $\mu=1,\ldots,M$ index activation samples, $i=1,\ldots,d$ activation
coordinates, and $j=1,\ldots,L$ learned SAE latents.  An observed activation
is $\bm h_\mu\in\mathbb R^d$, and the learned dictionary is
$D=[\bm d_1\ \cdots\ \bm d_L]\in\mathbb R^{d\times L}$, where
$\bm d_j=D_{:j}$ and $\lVert\bm d_j\rVert_2=1$.  Let
$s_{\mu j}\in\{0,1\}$ be the occupation of latent $j$ and $a_{\mu j}$ its
deterministic amplitude.  The random sparse code is
\begin{equation}
 z_{\mu j}=s_{\mu j}a_{\mu j},
 \qquad \bm z_\mu=\bm s_\mu\odot\bm a_\mu.
 \label{eq:latent-code}
\end{equation}
Conditioned on $\bm s_\mu$ and $\bm a_\mu$, the vector-output reconstruction
model is
\begin{align}
 \bm h_\mu
 &=\bm b+D\bm z_\mu+\bm\epsilon_\mu,
 &\bm\epsilon_\mu&\sim\mathcal N(\bm0,\beta^{-1}I_d),
 \label{eq:vector-reconstruction}\\
 h_{\mu i}
 &=b_i+\sum_{j=1}^{L}D_{ij}s_{\mu j}a_{\mu j}+\epsilon_{\mu i}.
 \label{eq:component-reconstruction}
\end{align}
Thus $D_{ij}$ is the decoder weight from latent $j$ to activation coordinate
$i$.  Relative to the scalar response $y^\mu$ in the original VG model, each
$s_{\mu j}$ here selects the entire vector atom $\bm d_j$.  The sample index is
also essential: the original VG selector is global, whereas an SAE has a
per-input occupation vector $\bm s_\mu$ and count $\lVert\bm z_\mu\rVert_0$.

The encoder supplies the mean-field selector posterior and the amplitudes,
\begin{align}
 Q_\phi(\bm s_\mu\mid\bm h_\mu)
   &= \prod_{j=1}^{L}\operatorname{Bernoulli}(s_{\mu j};m_{\mu j}),\\
 m_{\mu j}&=Q_\phi(s_{\mu j}=1\mid\bm h_\mu)
 =\sigma[g_{\phi,j}(\bm h_\mu-\bm b)],\\
 a_{\mu j}&=\operatorname{softplus}[r_{\phi,j}(\bm h_\mu-\bm b)],\\
 \overline z_{\mu j}
 &:=\E_{Q_\phi}[z_{\mu j}\mid\bm h_\mu]
 =m_{\mu j}a_{\mu j},\\
 \widehat{\bm h}_\mu
 &=\bm b+D\overline{\bm z}_\mu
 =\bm b+\sum_{j=1}^{L}m_{\mu j}a_{\mu j}\bm d_j .
 \label{eq:mean-reconstruction}
\end{align}
In particular, $m_{\mu j}a_{\mu j}$ is the posterior mean of $z_{\mu j}$,
not the binary-gated random code itself.  The expected per-sample
reconstruction energy is
\begin{align}
 E_\mu
 :={}&\E_{Q_\phi}\!\left[
 \frac12\left\lVert\bm h_\mu-\bm b-D(\bm s_\mu\odot\bm a_\mu)
 \right\rVert_2^2\right]\nonumber\\
 ={}& \frac12\left\lVert
 \bm h_\mu-\bm b-
 \sum_{j=1}^{L}m_{\mu j}a_{\mu j}\bm d_j
 \right\rVert_2^2 \nonumber\\
 &+\frac12\sum_{j=1}^{L}m_{\mu j}(1-m_{\mu j})a_{\mu j}^2
 \lVert\bm d_j\rVert_2^2 .
 \label{eq:energy}
\end{align}
The second term is not an auxiliary regularizer: it is the exact diagonal
Bernoulli variance correction under the mean-field posterior.

Following VG, we choose the factorized normalized sparse prior
\begin{align}
 P_\gamma(\bm s_\mu)
 &=\prod_{j=1}^{L}P_\gamma(s_{\mu j}),\nonumber\\
 P_\gamma(s_{\mu j})
 &=\frac{\exp(-\gamma s_{\mu j})}{1+\exp(-\gamma)},
 \qquad \gamma\in\mathbb R,
 \label{eq:prior}
\end{align}
so positive $\gamma$ suppresses occupation.  Write
\begin{align}
 H(m)&=-m\log m-(1-m)\log(1-m),\\
 R_\mu^{(\eta)}
 &=\gamma\sum_{j=1}^{L}m_{\mu j}
 +L\log(1+e^{-\gamma})
 -\eta\sum_{j=1}^{L}H(m_{\mu j}).
 \label{eq:selector-free-energy}
\end{align}
At the VG value $\eta=1$, $R_\mu^{(1)}$ is exactly
$D_{\mathrm{KL}}[Q_\phi(\bm s_\mu\mid\bm h_\mu)\Vert
P_\gamma(\bm s_\mu)]$; other entropy weights or removal of the entropy term
are ablations.  Equation~\eqref{eq:selector-free-energy}
uses the sign implied by Eq.~\eqref{eq:prior}: the printed signs of the field
term in Eqs.~(9) and (13) of Ref.~\cite{soh2025variational} are inconsistent
with that paper's normalized prior and its statement that increasing
$\gamma$ strengthens sparsity.

For a minibatch $\mathcal B\subseteq\{1,\ldots,M\}$ with
$B=|\mathcal B|$, the variational expectation is over the selectors only;
the deterministic amplitudes are point estimates produced by the encoder.
The Gaussian likelihood
contains $Bd$ scalar activation coordinates, whereas the selector prior and
entropy contain $BL$ binary variables.  In learned-precision
mode, $\theta_\beta=\log\beta$ is optimized and $\beta=e^{\theta_\beta}>0$.
The normalized minibatch-mean loss is
\begin{equation}
 \mathcal L_{\mathrm{learned}}
 =\frac1B\sum_{\mu\in\mathcal B}
 \left[
 e^{\theta_\beta}E_\mu
 -\frac d2\log\!\left(\frac{e^{\theta_\beta}}{2\pi}\right)
 +R_\mu^{(\eta)}
 \right].
 \label{eq:learned-beta-loss}
\end{equation}
The Gaussian normalization term in Eq.~\eqref{eq:learned-beta-loss} depends on
$\beta$ and therefore cannot be dropped.  For a fixed minibatch,
\begin{equation}
 \frac{\partial\mathcal L_{\mathrm{learned}}}{\partial\theta_\beta}
 =\frac{\beta}{B}\sum_{\mu\in\mathcal B}E_\mu-\frac d2,
 \qquad
 \beta_{\mathcal B}^\star
 =\frac{Bd}{2\sum_{\mu\in\mathcal B}E_\mu}.
 \label{eq:profiled-beta}
\end{equation}
The profiled-precision mode substitutes this minibatch optimum analytically.
The loss implemented in that mode is, up to a parameter-independent additive
constant,
\begin{widetext}
\begin{equation}
 \mathcal L_{\mathrm{profiled}}
 \doteq
 \frac1B\left[
 \frac{Bd}{2}\log\!\left(
 \frac{2\sum_{\mu\in\mathcal B}E_\mu}{Bd}
 \right)
 +\sum_{\mu\in\mathcal B}R_\mu^{(\eta)}
 \right].
 \label{eq:profiled-beta-loss}
\end{equation}
\end{widetext}
Here $\doteq$ omits $d[1+\log(2\pi)]/2$ per sample.  Thus the learned mode
optimizes one shared precision across minibatches, whereas the profiled mode
recomputes Eq.~\eqref{eq:profiled-beta} in every minibatch.  Neither mode fixes
$\beta$ externally; the implementation has no fixed-$\beta$ mode.

Because the present implementation has neither an amplitude posterior nor an
amplitude prior, Eqs.~\eqref{eq:learned-beta-loss} and
\eqref{eq:profiled-beta-loss} are not complete generative-model ELBOs; such a
model would add the corresponding amplitude KL term.

For a fitted model, the soft occupation observables are
\begin{align}
 \widehat L_0^{\mathrm{soft}}
 &=\frac1B\sum_{\mu\in\mathcal B}\sum_{j=1}^{L}m_{\mu j},\\
 \widehat\rho^{\mathrm{soft}}
 &=\frac{\widehat L_0^{\mathrm{soft}}}{L},\\
 V_{\mathrm{post}}&=\frac1{BL}
 \sum_{\mu\in\mathcal B}\sum_{j=1}^{L}m_{\mu j}(1-m_{\mu j}).
 \label{eq:observables}
\end{align}
We report these separately from the hard-thresholded firing count.  At
inference we use
\begin{align}
 \widehat s_{\mu j}&=\mathbf 1[m_{\mu j}>\tau],
 &z^{\mathrm{hard}}_{\mu j}&=\widehat s_{\mu j}a_{\mu j},\\
 \widehat L_0^{\mathrm{hard}}
 &:=\frac1B\sum_{\mu\in\mathcal B}\sum_{j=1}^{L}\widehat s_{\mu j},
 &\widehat\rho^{\mathrm{hard}}
 &:=\frac{\widehat L_0^{\mathrm{hard}}}{L}.
 \label{eq:hard-code}
\end{align}
The notation $L_0^\star$ is reserved for generative activity, whereas
$\widehat L_0^{\mathrm{soft}}$ is a VG posterior expectation and
$\widehat L_0^{\mathrm{hard}}$ is the measured SAE firing count.  Equivalently,
the latter is $B^{-1}\sum_{\mu\in\mathcal B}\sum_{j=1}^{L}
\mathbf 1[z^{\mathrm{hard}}_{\mu j}\ne0]$.  The
single-fit posterior quantity $V_{\mathrm{post}}$ is also distinct from the
cross-realization selection uncertainty studied by Soh \emph{et al.}; the
latter requires independent realizations.  A $\gamma$ sweep
therefore yields conditional soft and hard equations of state
$\widehat\rho^{\mathrm{soft}}(\gamma)$ and
$\widehat\rho^{\mathrm{hard}}(\gamma)$.

To make the transition estimator operational, at each target hard density we
align independently trained decoder atoms to a common $L$-slot reference.  We
use ground-truth atoms only in matched-width Stage 1 and a fixed reference SAE
otherwise.  On a common held-out batch, let
$\widetilde s^{(r)}_{\mu j}\in\{0,1\}$ be the aligned hard mask from repetition
$r=1,\ldots,R$ and define the SAE analogue of cross-realization selection
uncertainty,
\begin{align}
 \overline s_{\mu j}&=\frac1R\sum_{r=1}^{R}\widetilde s^{(r)}_{\mu j},\\
 \sigma_{\mathrm{sel}}^{\mathrm{SAE}}(\widehat\rho^{\mathrm{hard}})
 &=\frac1{BL}\sum_{\mu\in\mathcal B}\sum_{j=1}^{L}
 \overline s_{\mu j}(1-\overline s_{\mu j}).
 \label{eq:sae-selection-uncertainty}
\end{align}
For compactness, write $\rho=\widehat\rho^{\mathrm{hard}}$ in the following
template.  We fit the full nondegenerate density curve, rather than selecting its raw
global minimum, to the mean-field family of Ref.~\cite{soh2025variational},
\begin{equation}
 f(\rho;\rho_\star)=
 \begin{cases}
 \dfrac{\rho}{\rho_\star}(\rho_\star-\rho),&\rho\leq\rho_\star,\\[4pt]
 (\rho-\rho_\star)\dfrac{1-\rho}{1-\rho_\star},&\rho>\rho_\star.
 \end{cases}
 \label{eq:mean-field-template}
\end{equation}
A nonnegative mixture over candidate $\rho_\star$ values gives a normalized
weight distribution and hence a point estimate and interval.  This transfer
from global VG variables to sample-dependent SAE gates is a hypothesis to be
calibrated in Stage 1, not an assumed identity.  The equation-of-state slope
and $c_{\mathrm{dec}}$ remain complementary cross-checks.

Unlike the original global VG precision, the profiled mode is a
log-batch-energy surrogate and is unbounded as the residual energy approaches
zero.  Confirmatory noiseless and overcomplete experiments therefore compare
it with learned precision.

\section{Three-stage experimental program}
\label{sec:experiments}

All architectures are compared at a common measured test-set
$\widehat L_0^{\mathrm{hard}}$, not at nominal penalty values; VG soft
occupation is shown separately.
Reconstruction is retained as a fidelity check, while feature recovery or
downstream behavior is the primary endpoint.

\subsection{Stage 1: finalized controlled baseline}

The controlled baseline has ground-truth width $L^\star=128$ and learned width
$L=128$.  It uses the matched-capacity sparse-coding model
\begin{align}
 s^\star_{\mu k}&\sim\operatorname{Bernoulli}(p_k^\star),\\
 a^\star_{\mu k}&\sim\operatorname{Exponential}(1),\\
 z^\star_{\mu k}&=s^\star_{\mu k}a^\star_{\mu k},\\
 \bm h_\mu&=D^\star\bm z^\star_\mu,
 \label{eq:stage1-data}
\end{align}
where $k=1,\ldots,L^\star$ indexes ground-truth latents and
\begin{equation}
 p_k^\star=0.1\,
 \frac{k^{-1/2}}{128^{-1}\sum_{\ell=1}^{128}\ell^{-1/2}}.
\end{equation}
The columns of $D^\star\in\mathbb R^{16\times128}$ are independent Gaussian
directions normalized to unit length.  ``No added coherence'' means that no
common direction is injected; it does not mean orthogonality, which is
impossible for 128 atoms in 16 dimensions.  Supports are independent, and the
model has no hierarchy, correlated firing, or observation noise.  The train
and test sets each contain 512 samples.  The ground-truth and learned widths
are both 128, so
\begin{equation}
 L_0^\star=\sum_{k=1}^{L^\star}p_k^\star=12.8,
 \qquad \rho^\star=0.1.
\end{equation}
This is an expected count, not a fixed per-sample cardinality.

Each completed precision-mode calibration uses seed 1, 1000 optimization steps, batch size
128, learning rate $10^{-2}$, gradient-norm clipping at 1, and the final
(\texttt{last}) checkpoint; metrics are recorded every 25 steps.  The dead
feature window is 100 steps, the evaluation dead threshold is $10^{-6}$, and
the VG hard-mask threshold is $0.5$.  The sweep contains 273 runs: 33 VG-SAE, 16
L1/ReLU, 128 TopK, 41 BatchTopK, 32 JumpReLU, and 23 Gated SAE control values.
Within each architecture, control values share the data, initialization seed,
and batch order.  The five baselines use the pinned SAELens v6.47 training
implementations.  VG-SAE uses its Bernoulli variance term, unit entropy
coefficient.  VG-SAE is evaluated in separate profiled- and learned-precision
sweeps; the stored $\beta=1$ is an initialization/reporting value, not a fixed
precision.  Each single-seed sweep is a calibration, not an
uncertainty estimate.  In the confirmatory study, each data seed fixes the
dictionary and held-out inputs across optimization repetitions; curves are
estimated after decoder alignment and then aggregated across data seeds on a
common $\widehat L_0^{\mathrm{hard}}$ grid.  The stored calibration field
\texttt{rho\_model} is soft for VG-SAE but hard for the five baselines, so it is
not by itself the final cross-method axis.

Because the full generative state is known, the existing sweep records decoder
matching, support precision/recall/F1, feature coverage, and amplitude
shrinkage alongside MSE and explained variance.  The confirmatory extension
will add multi-seed stability and $c_{\mathrm{dec}}$ relative to both the true
dictionary and the finite-dimensional random-direction baseline
\begin{equation}
 c_{\mathrm{null}}(d)=
 \frac{\Gamma(d/2)}{\sqrt{\pi}\,\Gamma[(d+1)/2]},
 \qquad c_{\mathrm{null}}(16)\simeq0.203.
\end{equation}
This baseline is required because the orthogonality assumptions underlying the
toy analysis of Ref.~\cite{chanin2026sparse} do not hold for 128
atoms in 16 dimensions.  The key falsifiable prediction is that recovery is
optimized near $L_0^\star=12.8$ even if a different occupation reconstructs
better.  We will then test whether the VG transition and
Eq.~\eqref{eq:cdec} locate that regime.

\subsection{Stage 2: official SynthSAEBench}

The second stage will use the
\href{https://huggingface.co/decoderesearch/synth-sae-bench-16k-v1}{released
SynthSAEBench-16k generator} at the pinned
\href{https://huggingface.co/decoderesearch/synth-sae-bench-16k-v1/tree/b2efd8b919ae46d6d487c73d46db5ee52813621d}{\texttt{b2efd8b}}
revision rather than regenerating an approximately equivalent model.  It has
$d=768$, ground-truth width $L^\star=16{,}384$, learned SAE width $L=4096$,
and mean full-generator
$L_0^\star$ of approximately 34.  Its skewed firing rates, hierarchy, mutual
exclusion, low-rank support correlations, heterogeneous amplitudes, and bias
provide a structured test beyond Stage 1.

The full protocol uses the official pretrained generator, batch size 1024,
199,999,488 online training samples, and 24,999,936 independently seeded
evaluation samples.  VG-SAE is run separately with profiled and learned
precision, with mode-specific sparsity calibration.
The same six architectures will be evaluated with native SynthSAEBench metrics:
hard SAE $\widehat L_0^{\mathrm{hard}}$ and true $L_0^\star$,
feature-classification
precision/recall/F1, Matthews
correlation, uniqueness, explained variance, dead latents, and shrinkage.  VG
posterior observables will remain separate diagnostics; frequency-stratified recovery,
duplicate/missed features, and $c_{\mathrm{dec}}$ are planned additions.

Here the ground-truth vocabulary is four times wider than the SAE\@.  Full
dictionary recovery is therefore impossible, and the generator's full
$L_0^\star$ divided by SAE width is not a uniquely correct SAE density.  We instead
will measure recovery by feature frequency, duplicate and missed features, and
capacity-limited coverage.  Custom hierarchy or correlation ablations will be
reported separately from the fixed official benchmark.

\subsection{Stage 3: real activation data}

The third stage will use versioned, train/validation-separated language-model
activation stores through the same SAELens interface.  The current scaffold is
a VG-only sweep on at most 4096 WikiText tokens from the GPT-2-small layer-8
residual stream; it is not the confirmatory six-method experiment.  The final
model, layer, corpus, and token budget remain to be fixed and will be recorded
before running the sweep.  That experiment will fit normalization on training
data only and match width, tokens, optimizer budget, and measured
$\widehat L_0^{\mathrm{hard}}$ across methods.

There is no observed true dictionary or true $L_0$ in this stage.  We therefore
call the output an inferred operating density, not the actual number of LLM
features.  Candidate densities will be chosen without using reconstruction as
the sole criterion, then evaluated by held-out explained variance and loss
recovery, dead and firing-frequency distributions, seed stability,
$c_{\mathrm{dec}}$, sparse probes, and causal ablation or patching.  Agreement
between the VG transition, decoder diagnostics, and downstream optima would
support the order-parameter interpretation; systematic disagreement would
falsify the proposed estimator.

\section{Claim boundary and outlook}

The intended take-home message is not that any physical metaphor makes
$L_0$ identifiable.  It is that reconstruction error does not define the
feature decomposition, and that $L_0$ should be treated as an inferential
target under an explicit probabilistic model rather than chosen as an
arbitrary budget.  In the current
VG-SAE, $\gamma$ remains an external field, so the method is not yet
parameter-free.  After specifying a proper generative model for amplitudes,
evidence optimization or a hyperprior over the inclusion rate is a natural
next step.  Until then, the paper's target empirical claim
is the testable co-localization of free-energy transition diagnostics and the
feature-recovery regime, with exact validation restricted to synthetic data.

\bibliography{apssamp}

\end{document}
